\uuid{TJ6f}
\titre{Structure de  sous-espace vectoriel}
\auteur{hanani ; mzari}
\organisation{exo7}
\licence{CC BY-NC-SA 4.0}
\datecreate{2020-01-20}
\niveau{L1}
\module{Algèbre}
\chapitre{Espace vectoriel}
\sousChapitre{Définition, sous-espace}
\difficulte{2}
\theme{Espaces vectoriels, Structure de sous-espace vectoriel}
\exo7qcmid{lille-2.espaces-vectoriels.010}

\contenu{
\begin{question}
\qclassification{106.01}
\qtags{motcle=Structure de  sous-espace vectoriel}
Soit $E$ un espace vectoriel. Quelles sont les assertions vraies ?
\begin{answers}
\bad{L'intersection de deux sous-espaces vectoriels de $E$ peut être vide.}
\bad{Si $F$ est un sous-espace vectoriel de $E$, alors $F$ contient toute combinaison linéaire d'éléments de $E$.}
\good{Il existe un sous-espace vectoriel de $E$ qui contient un seul élément.}
\good{Si $F$ est un sous-ensemble non vide de $E$ qui contient  toute combinaison linéaire de deux vecteurs de $F$, alors $F$ est un sous-espace vectoriel de $E$.}
\end{answers}
\begin{explanations} L'intersection de deux sous-espaces vectoriels de $E$ contient au moins le vecteur nul.
Le seul sous-espace vectoriel de $E$ qui contient un seul élément est $\{0_E\}$, où $0_E$ est le zéro de $E$.\\
Un sous-ensemble non vide de $E$ est un sous-espace vectoriel de $E$ si et seulement
s'il contient toute combinaison linéaire d'éléments de $F$. Ceci revient à dire que $F$ contient
toute combinaison linéaire de deux éléments de $F$.
\end{explanations}
\end{question}
}
